\documentclass[11pt]{article}

\usepackage[T1]{fontenc}
\usepackage{lmodern}
\usepackage{amsmath,amssymb,amsthm,mathtools}
\usepackage{booktabs,array}
\usepackage{enumitem}
\usepackage[a4paper,margin=27mm]{geometry}
\usepackage[colorlinks=true,linkcolor=blue,citecolor=blue,urlcolor=blue]{hyperref}
\setlength{\emergencystretch}{2em}
\numberwithin{equation}{section}

\newtheorem{theorem}{Theorem}[section]
\newtheorem{proposition}[theorem]{Proposition}
\newtheorem{lemma}[theorem]{Lemma}
\newtheorem{corollary}[theorem]{Corollary}
\theoremstyle{definition}
\newtheorem{example}[theorem]{Example}
\theoremstyle{remark}
\newtheorem{remark}[theorem]{Remark}

\newcommand{\CC}{\mathbf C}
\newcommand{\ZZ}{\mathbf Z}
\newcommand{\RR}{\mathbf R}
\newcommand{\PP}{\mathbf P}
\newcommand{\one}{\mathbf 1}
\newcommand{\Reg}{\operatorname{Reg}}
\newcommand{\Aug}{\operatorname{Aug}}
\newcommand{\mob}{\mu_{\mathrm{Mob}}}
\DeclareMathOperator{\ord}{ord}
\DeclareMathOperator{\Res}{Res}
\DeclareMathOperator{\Tr}{Tr}
\DeclareMathOperator{\relint}{relint}
\DeclareMathOperator{\mult}{mult}

\title{Contact Reconstruction and Spectral Balance\\
for Surface Singularities with a Multiple Smooth Tangent Cone}
\author{YINGFENG JIANG}
\date{19 October 2026}

\begin{document}
\maketitle

\begin{abstract}
We recover the degrees and all contact orders of surface germs
$f=g_e^r+h_{er+1}$, $e\ge3$, $r\ge2$, where $V(g_e)\subset\PP^2$
is smooth and $V(h_{er+1})$ is smooth at its finite intersection
with $V(g_e)$. The monodromy characteristic polynomial yields
these data by coefficient extraction and finite M\"obius inversion;
the local topological zeta function yields them through its
poles and residues. Thus either invariant determines the other
and classifies semisimplified Milnor monodromy in this family.
We also obtain the spectral balance
$\mu+M_r-M_1=(er)^3-(er)^2+er$, where $M_m$ is the multiplicity
of an individual primitive $m$th root on $H^2$. At fixed $e,r$,
it forces every monodromy-equivariant injection to be an
isomorphism. We realize distinct contact profiles with equal
Milnor numbers and incomparable monodromy, and describe the
spectral change in a contact collision. Explicit formulas for
the two invariants also yield simple actual poles, eigenvalues
at the origin, and twisted topological holomorphy.
\end{abstract}

\medskip
\noindent\textbf{Keywords.}
Surface singularities; monodromy; topological zeta functions;
tangency; contact reconstruction.

\section{Introduction}\label{sec:introduction}

We study how finite tangency affects the monodromy of surface
germs with a multiple smooth tangent cone. The main result
recovers both degrees and all contact orders from either the
monodromy characteristic polynomial or the local topological
zeta function. A spectral balance then governs comparisons
between the monodromy representations of different profiles.

We work over $\CC$. Let $g_e,h_D\in\CC[x,y,z]$ be nonzero homogeneous
polynomials of degrees $e,D$, respectively, and set
\begin{equation}\label{eq:family}
 f=g_e^r+h_D,\qquad e\ge3,\quad r\ge2,\quad A=er,\quad D=A+1.
\end{equation}
Assume that the zero scheme $C=V(g_e)\subset\PP^2$ is smooth,
that $\mathcal P=C\cap H$ is finite for $H=V(h_D)$, and that
$H$ is smooth at every point of $\mathcal P$.
For $p\in\mathcal P$, put
$q_p=I_p(C,H)$. The finite multiset $\{q_p\}$ is the
\emph{contact profile}. These germs have isolated critical points,
as proved in Lemma~\ref{lem:isolated}.

Write $Z_f(s)=Z_{\mathrm{top},0}(f;s)$ for the local topological
zeta function with standard volume form. Let $F_f$ be the Milnor
fiber at the origin, $T_f$ its monodromy on $H^2(F_f,\CC)$, and
\[
 \Delta_f(t)=\det(1-tT_f\mid H^2(F_f,\CC)).
\]
Write $\mu_f$ for its Milnor number.
Our first result reconstructs the entire profile from either
one of these invariants.

\begin{theorem}[Reconstruction]\label{thm:main}
Each of $\Delta_f$ and $Z_f$ uniquely determines
$(e,r,\{q_p\})$ in the family \eqref{eq:family}.
For two members, equality of these data is equivalent to
equality of either invariant and to isomorphism of their
semisimplified Milnor monodromy. Recovery from $\Delta_f$ uses
coefficient extraction and finite M\"obius inversion.
\end{theorem}

For $\Delta_f$, the initial step is visible in a single term. If
$P_f=(1-t)\Delta_f$ begins as $1-\kappa t^A+\cdots$, then
\begin{equation}\label{eq:initial-recovery-intro}
 A=er,\qquad \kappa=e^2-3e+3,\qquad
 e=\frac{3+\sqrt{4\kappa-3}}2,\qquad r=\frac Ae.
\end{equation}
Two further coefficients recover the number of intersection
points and the number of contacts of order at least two.
The logarithmic derivative of the remaining contact polynomial
then recovers each order and its frequency.
For $Z_f$, the largest actual pole and its residue first
recover $e,r$; the remaining contact poles and their residues
recover the higher contact orders and their frequencies;
B\'ezout gives the number of transverse intersections
(Proposition~\ref{prop:z-recovery}). The recovered data also
determine every twisted topological zeta function
(Corollary~\ref{cor:twisted-recovery}).

A collision of $q$ transverse intersections into one contact of
order $q$ increases the Milnor number by $q-1$. At the same time
it decreases the multiplicity of every primitive $r$th root of
unity by $q-1$. Our second result expresses this compensation
as a balance valid throughout the family.
Let $M_m(f)$ denote the algebraic multiplicity of an individual primitive
$m$th root as an eigenvalue on $H^2(F_f,\CC)$.

\begin{theorem}[Spectral balance and incomparability]\label{thm:balance}
For every member of \eqref{eq:family},
\begin{equation}\label{eq:balance-intro}
 \mu_f+M_r(f)-M_1(f)=A^3-A^2+A.
\end{equation}
For two members with the same $e,r$, every monodromy-equivariant injection
$H^2(F_f,\CC)\hookrightarrow H^2(F_{\widetilde f},\CC)$ is an
isomorphism. If their contact profiles differ, neither
direction admits such an injection, also after semisimplification.
\end{theorem}

This incomparability occurs even at equal Milnor number.
For $e=4,r=2$, the profiles
$\{2,4,1^{30}\}$ and $\{3,3,1^{30}\}$ both occur and both give
$\mu=419$. The first has primitive $20$th eigenvalues and the
second primitive $30$th eigenvalues, with neither order
present in the other
(Example~\ref{ex:equal-mu}).

Both reconstruction procedures start from the same calculation.
Blowing up the origin gives three local models, over
$\PP^2\setminus C$, $C\setminus\mathcal P$, and
$\mathcal P$. Their weighted Newton and monodromy formulas
assemble into a polynomial contact transform for $\Delta_f$
and a rational one for $Z_f$.

The resolution formula for the topological zeta function is due
to Denef and Loeser~\cite{DL92}, and the monodromy computation
uses A'Campo's formula~\cite{AC75}. The partial-resolution
formula of Gusein-Zade, Luengo, and Melle-Hern\'andez
\cite[Theorem~1 and Corollary~1]{GLM97} provides the
multiplicative assembly; their \S3.5 includes the transverse
specialization of our monodromy formula. Melle-Hern\'andez's
formula, recorded in~\cite[\S3.4, Remark~2]{GLM97}, gives the
Milnor number for arbitrary contacts. Building on these
forward calculations, we recover the full contact data
from either invariant and derive the spectral comparison.

Related monodromy results concern Newton nondegenerate surfaces
\cite{LVP11}, their local motivic zeta functions
\cite[Theorem~7.2.1]{LPS26}, and generalized suspensions and
L\^e--Yomdin surfaces~\cite{AGGL26}.
The present tangent cone is the non-reduced divisor $rC$;
the surface L\^e--Yomdin class has a reduced tangent cone
\cite[Definition~7.1]{AGGL26}. Likewise, the lowest-degree
face polynomial $g_e^r$ and all its first derivatives vanish
on the nonempty torus part of $C$, while the three models after
blowing up the origin are Newton nondegenerate. For a broader
account of the monodromy conjecture, see~\cite{Veys24}.

Sections~\ref{sec:geometry}--\ref{sec:global} compute the local
models and global invariants. Section~\ref{sec:reconstruction}
inverts the two contact transforms, placing the pole calculation
alongside the reconstruction from $Z_f$.
Section~\ref{sec:collisions} proves spectral balance and treats
contact collisions and extremal members of the family.
Section~\ref{sec:twisted} gives the twisted formulas and
holomorphy. The appendices contain the cone calculations
and a worked example of reconstruction.

\section{Three local models after blowing up the origin}
\label{sec:geometry}

\subsection{Contact data}

Write
\begin{equation}\label{eq:contact-data}
 \begin{gathered}
 c=\chi(C)=e(3-e),\qquad \kappa=3-c,\qquad
 n=eD=\sum_{p\in\mathcal P}q_p,\qquad b=|\mathcal P|,\\
 n_q=|\{p:q_p=q\}|,\qquad
 S=\sum_{q\ge2}(q-1)n_q=n-b,\\
 B=(r-1)D+1,\qquad \beta=3r-2.
 \end{gathered}
\end{equation}
Here $\chi$ denotes Euler characteristic with compact supports;
$c\le0$ and $b\ge1$.

\subsection{Resolution conventions and isolatedness}

Monodromy is taken along a counterclockwise loop in the base.
For a log resolution $h:Y\to\CC^3$, let $E_i$ be the components
of the total transform of $f^{-1}(0)$ and write
\[
 \operatorname{div}(f\circ h)=\sum_iN_iE_i,\qquad
 \operatorname{div}(h^*(dx\wedge dy\wedge dz))
       =\sum_i(\nu_i-1)E_i.
\]
For a nonempty index set $I$, put
$E_I^\circ=(\bigcap_{i\in I}E_i)\setminus\bigcup_{j\notin I}E_j$.
Our normalization is
\begin{equation}\label{eq:z-definition}
 Z_f(s)=\sum_{\varnothing\ne I}
  \chi(E_I^\circ\cap h^{-1}(0))
  \prod_{i\in I}(N_i s+\nu_i)^{-1}.
\end{equation}
For a local function $F$, let $\omega$ be a nonzero
holomorphic top-degree form with
$\operatorname{Supp}(\operatorname{div}\omega)\subset F^{-1}(0)$.
On a common log resolution of $F$ and $\omega$, taken to be an
isomorphism outside $F^{-1}(0)$, use \eqref{eq:z-definition}
with $\nu_i-1=\ord_{E_i}(h^*\omega)$.
The form may vanish along $F^{-1}(0)$.

Auxiliary coordinate boundaries with data $(N,\nu)=(0,1)$
may be used to refine the Euler strata. They are additional
to the components of the function divisor: their denominator
factor is $1$, so Euler additivity preserves the sum.

For any local Milnor fiber we take
\begin{equation}\label{eq:zeta-convention}
 \zeta(t)=\prod_j\det(1-tT_j\mid H^j)^{(-1)^{j+1}}.
\end{equation}
This is the reciprocal of the convention in~\cite{GLM97}.
In this convention A'Campo's formula reads
\begin{equation}\label{eq:acampo}
 \zeta(t)=\prod_{i:N_i>0}(1-t^{N_i})^{
       -\chi(E_i^\circ\cap h^{-1}(0))}.
\end{equation}
In \eqref{eq:acampo}, $E_i^\circ$ removes all other positive
multiplicity components, including the strict transform.

\begin{lemma}\label{lem:isolated}
The function $f$ in \eqref{eq:family} is reduced and has an
isolated critical point at the origin. Its Milnor fiber is
connected, with reduced cohomology concentrated in degree $2$.
In particular,
\begin{equation}\label{eq:zeta-delta}
 \zeta_f(t)=\bigl((1-t)\Delta_f(t)\bigr)^{-1}.
\end{equation}
\end{lemma}

\begin{proof}
At a nonzero point with $f=df=0$, the homogeneous Euler identity
and the equation $f=0$ give
\[
 A g_e^r+D h_D=0,\qquad g_e^r+h_D=0.
\]
Since $D-A=1$, both $g_e$ and $h_D$ vanish. The corresponding
projective point lies in $\mathcal P$, where $dh_D\ne0$.
But $r\ge2$ gives $df=dh_D$ at that point, a contradiction.

If the critical locus had a positive dimensional germ through
the origin, it would contain an analytic curve through the
origin. The restriction of $f$ to that curve would be constant,
hence zero, contrary to what was just proved. A repeated
factor of $f$ would also give a positive dimensional critical
locus. The Milnor bouquet theorem~\cite{Mil68} now gives the
cohomology statement and \eqref{eq:zeta-delta}.
\end{proof}

\subsection{The local equations}

Let $\pi_0:X_1\to\CC^3$ be the blowup of the origin and let
$E=\pi_0^{-1}(0)\simeq\PP^2$. Near $C$ choose $u$ normal to $E$,
$v$ an equation for $C$ in $E$, and $w$ along $C$. Then
\begin{equation}\label{eq:blowup-equation}
 f\circ\pi_0=u^A(v^r+u h(v,w)),\qquad
 \pi_0^*(dx\wedge dy\wedge dz)
   =u^2\cdot\text{unit}\,du\wedge dv\wedge dw.
\end{equation}
Outside $C$ the function is $u^A$ times a unit. On
$C\setminus\mathcal P$, replacing $u$ by $h(v,w)u$ gives
\[
 G=u^A(v^r+u).
\]
At a point of contact order $q\ge2$, smoothness of $H$ implies
$\partial_v h(0)\ne0$. The implicit function theorem gives
$h=\alpha(v,w)(v+\phi(w))$, where $\alpha$ is a unit and
$\ord_0\phi=q$. A change of the coordinate $w$ makes
$\phi(w)=w^q$, and the change $u'=\alpha u$ gives
\begin{equation}\label{eq:tangency-model}
 \mathcal T_q=u^A(v^r+uv+uw^q).
\end{equation}
When $q=1$, take $w'=h-v$ to obtain the same expression.
These changes multiply the function and the top-degree form by
units. Their orders and divisor strata on a common resolution
are therefore unchanged. The three local models are computed
with the form
\[
 \omega=u^2\,du\wedge dv\wedge dw.
\]

\subsection{Euler assembly}

The three local computations assemble by Euler integration.
For monodromy this is the partial-resolution formula of
\cite[Theorem~1 and Corollary~1]{GLM97}.

\begin{lemma}\label{lem:assembly}
Let $Z_G,Z_q$ be the local topological functions of $G,\mathcal T_q$
with the form $\omega$, and let $\zeta_G,\zeta_q$ be their
monodromy zeta functions. Then
\begin{align}
 Z_f&=\kappa Z_{u^A}+(c-b)Z_G+\sum_{q\ge1}n_qZ_q,
       \label{eq:assembly-z}\\
 \zeta_f&=\zeta_{u^A}^{\,\kappa}
            \zeta_G^{\,c-b}\prod_{q\ge1}\zeta_q^{\,n_q}.
       \label{eq:assembly-zeta}
\end{align}
\end{lemma}

\begin{proof}
Choose a log resolution $\rho:Y\to X_1$ of the pulled-back
function and form, and set $h=\pi_0\rho$.
The pulled-back form on $Y$ is $h^*(dx\wedge dy\wedge dz)$,
so its numerical data coincide with those of the original germ.
For each resolution stratum, Euler Fubini gives
\[
 \chi(E_I^\circ\cap h^{-1}(0))
   =\int_E\chi(E_I^\circ\cap\rho^{-1}(y))\,d\chi(y).
\]
Multiplying by its fixed denominator and summing yields the
Euler integral of the local weighted topological function.
Local changes of variables preserve it by
\cite[Proposition~1.3]{AGGL26}. The corresponding normal
crossings identity will also be used for the twisted functions.
Blowing up a smooth center of codimension $v$
having normal crossings with the boundary and contained in
exactly $k$ of its components replaces their
linear denominators $L_i$ by an additional denominator
$L_*=\sum_{i=1}^kL_i+v-k$. The nonzero Euler strata of the
projective exceptional fiber give
\begin{equation}\label{eq:snc-identity}
 \frac1{L_*}\left[
   \sum_{i=1}^k\prod_{j\ne i}L_j^{-1}
       +(v-k)\prod_{i=1}^kL_i^{-1}\right]
       =\prod_{i=1}^kL_i^{-1}.
\end{equation}
Boundary components transverse to the center contribute
common factors on the two sides. The identity therefore
applies with the orders of the pulled-back form.

For monodromy, on this fixed resolution put
\[
 a_N(y)=\sum_{i:N_i=N}
       \chi(E_i^\circ\cap\rho^{-1}(y)).
\]
The local A'Campo expression is
$\zeta_{f\circ\pi_0,y}=\prod_N(1-t^N)^{-a_N(y)}$.
Euler Fubini identifies the global exponent with
$-\int_E a_N\,d\chi$ for each $N$, giving the multiplicative
Euler integral of these local expressions. Finally stratify
\[
 E=(\PP^2\setminus C)\sqcup(C\setminus\mathcal P)
      \sqcup\bigsqcup_{p\in\mathcal P}\{p\},
\]
whose Euler characteristics are $\kappa,c-b,1$. The local
resolution invariants are constant on these strata, giving
\eqref{eq:assembly-z} and \eqref{eq:assembly-zeta}.
\end{proof}

\section{The local tangency calculation}\label{sec:local}

Use the linear functions
\begin{equation}\label{eq:linear-functions}
 a=As+3,\quad Q=rDs+3r+1,\quad d=s+1,\quad
 J=Ds+4,\quad K_q=qBs+q\beta+1.
\end{equation}
Three identities used below are
\begin{equation}\label{eq:linear-identities}
 Q=ra+rd-(r-1)=rJ-(r-1),\qquad
 K_q=q(Q-a)+1,\qquad B-(r-1)D=1.
\end{equation}

\begin{proposition}[Local contact formula]\label{prop:local}
For the local models in Section~\ref{sec:geometry},
\begin{align}
 Z_{u^A}&=\frac1a,&
 \zeta_{u^A}&=(1-t^A)^{-1},\label{eq:local-monomial}\\
 Z_G&=\frac{r-1+a^{-1}+d^{-1}}Q,&
 \zeta_G&=\frac{1-t^{rD}}{1-t^D},\label{eq:local-g}\\
 Z_q&=\frac{a^{-1}+r/d}{Q}
       -\frac{(q-1)rsJ}{dQK_q},&
 \zeta_q&=\frac{1-t^B}{1-t^{qB}}.\label{eq:local-q}
\end{align}
\end{proposition}

\begin{proof}
\emph{The topological function of $\mathcal T_q$.}
The Newton vertices of $\mathcal T_q$ are
\[
 X=(A,r,0),\qquad Y=(D,1,0),\qquad Z=(D,0,q).
\]
The primitive facet normals and their numerical data for
$\omega=u^2du\wedge dv\wedge dw$ are
\begin{center}
\begin{tabular}{lrrl}
\toprule
Normal & $N$ & $\nu$ & Denominator\\
\midrule
$\mathsf U=(1,0,0)$ & $A$ & $3$ & $a$\\
$\mathsf V=(0,1,0)$ & $0$ & $1$ & $1$\\
$\mathsf W=(0,0,1)$ & $0$ & $1$ & $1$\\
$\mathsf R=(r,1,0)$ & $rD$ & $3r+1$ & $Q$\\
$\mathsf S=(0,q,1)$ & $q$ & $q+1$ & $p_q=qs+q+1$\\
$\mathsf T=(q(r-1),q,1)$ & $qB$ & $q\beta+1$ & $K_q$\\
\bottomrule
\end{tabular}
\end{center}
All compact edges have normalized length one. The cross product
of $Y-X$ and $Z-X$ is $-\mathsf T$, so the compact triangle has
normalized area one. The edge polynomials are primitive
binomials, and the triangle polynomial has nonzero $w$-derivative
on the torus. They are nondegenerate. A regular subdivision of
the Newton fan consequently gives a toric log resolution near
the fiber over the origin.

The weighted orbit calculation in Appendix~\ref{app:cones}
gives
\begin{equation}\label{eq:local-raw}
 Z_q=\frac1{K_q}\left[
  \frac q a+\frac1{aQ}+\frac{r-1}Q-(q-1)
          +\frac{qr-1+Q^{-1}}d\right].
\end{equation}
Using \eqref{eq:linear-identities}, its numerator satisfies
\[
 adQK_q Z_q=K_q(d+ra)-(q-1)ra(d-1)(a+d).
\]
Since $d-1=s$ and $a+d=J$, this proves the first formula
in \eqref{eq:local-q}.

\emph{The monodromy of $\mathcal T_q$.}
In the cone
dual to $XY$, insert the rays
\[
 Q_j=(j(r-1),j,1),\qquad 0\le j\le q.
\]
Here $Q_0=\mathsf W$, $Q_q=\mathsf T$, and the cross product
of consecutive vectors is $(-1,r-1,0)$. These two-dimensional
cones are regular and can be kept in a regular refinement of
the full fan.

In A'Campo's formula, the two-dimensional torus orbit of
$E_{\mathsf T}$ has complement of Euler characteristic $+1$:
the strict transform is $1+\xi+\eta=0$ in $(\CC^*)^2$.
An orbit of a two-dimensional cone can remain in a single
positive multiplicity divisor only if its other ray has
$N=0$. This forces $\rho_u=0$ and $\rho_v\rho_w=0$,
so the only such rays are $\mathsf V,\mathsf W$.
The cones containing $\mathsf V$ give Euler characteristic
zero. Among the cones containing $\mathsf W$, the sole
nonzero contribution comes from the edge $XY$. It is
$\CC^*$ minus one point, has Euler characteristic $-1$,
and belongs to the divisor $E_{Q_1}$, of multiplicity $B$.

There is no three-dimensional cone with just one positive
multiplicity ray: that would require both $\mathsf V$ and
$\mathsf W$, which lie in no common Newton cone. All other
orbits have an extra torus factor and Euler characteristic
zero. Also, the strict transform's own open stratum misses
the fiber, since every orbit over the origin lies in a
positive multiplicity toric divisor. Thus the divisor Euler
numbers are $+1$ at $\mathsf T$ and $-1$ at $Q_1$. For $q=1$
these are the same divisor and the contributions cancel.
These cases exhaust the fiber orbits relevant to A'Campo's
formula in every regular refinement preserving the indicated
two-dimensional cones.
Formula \eqref{eq:acampo} now gives $\zeta_q$.

\emph{The other two models.}
For $G$, take a plane resolution in $u,v$ and its product
with the smooth $w$-direction. The Newton vertices are
$(A,r)$ and $(D,0)$, with normals $(1,0),(r,1),(0,1)$
and denominators $a,Q,1$. The two vertex cones have lattice
indices $1,r$, giving $1/(aQ)$ and $r/Q$; the primitive
edge gives $-s/(dQ)$ by the orbit computation.
Their sum is \eqref{eq:local-g}.
For its monodromy insert the rays $(j,1)$, $1\le j\le r$,
between $(0,1)$ and $(r,1)$. The divisor at $(r,1)$ has open
Euler characteristic $-1$, and the divisor at $(1,1)$ has
an additional point of Euler characteristic $+1$ next to
the zero multiplicity boundary. Their multiplicities are
$rD$ and $D$. This gives $\zeta_G$. The monomial formulas
follow immediately.
\end{proof}

\section{Contact transforms and global invariants}\label{sec:global}

\subsection{The two contact transforms}

The local formulas express the global invariants through the
contact polynomial and the rational contact transform:
\begin{equation}\label{eq:contact-transforms}
 \Gamma_f(u)=\prod_{q\ge2}(1+u+\cdots+u^{q-1})^{n_q},
 \qquad
 U_f(x)=\sum_{q\ge2}\frac{n_q(q-1)}{1+qx}.
\end{equation}
The total contact excess is
$\deg\Gamma_f=U_f(0)=S$.

\begin{theorem}[Contact formulas]\label{thm:formulas}
The global invariants satisfy
\begin{align}
 Z_f(s)={}&\frac{\kappa}{a}
       +\frac{(c-b)(r-1)}Q+\frac{c}{aQ}
       +\frac{c+b(r-1)}{dQ}
       -\frac{rsJ}{dQ}U_f(Bs+\beta),\label{eq:global-z}\\
 \Delta_f(t)={}&\frac{(1-t^A)^\kappa}{1-t}
       \left(\frac{1-t^{rD}}{1-t^D}\right)^{b-c}
       \Gamma_f(t^B).\label{eq:global-delta}
\end{align}
In particular,
\begin{equation}\label{eq:global-mu}
 \begin{aligned}
 \mu_f&=A\kappa+(r-1)D(b-c)+B(n-b)-1\\
      &=A^3-A^2+A-1+c-b.
 \end{aligned}
\end{equation}
\end{theorem}

\begin{proof}
Insert Proposition~\ref{prop:local} into
Lemma~\ref{lem:assembly}. For the topological function, the
coefficients of $1/(aQ)$ combine to $c$, and those of $1/(dQ)$
combine to $c+b(r-1)$. The remaining contact sum is exactly
$U_f(Bs+\beta)$. For monodromy, \eqref{eq:assembly-zeta} gives
\[
 \zeta_f=(1-t^A)^{-\kappa}
  \left(\frac{1-t^{rD}}{1-t^D}\right)^{c-b}
  \prod_{q\ge2}
       \left(\frac{1-t^B}{1-t^{qB}}\right)^{n_q}.
\]
Use \eqref{eq:zeta-delta} to obtain \eqref{eq:global-delta}.
Its degree is \eqref{eq:global-mu}. Polynomiality is also
visible directly: the geometric-series quotients are
polynomials, and $(1-t^A)^\kappa/(1-t)$ is a polynomial.
Using $n=eD$ and $B-(r-1)D=1$ simplifies the degree formula
to its second expression.
\end{proof}

The classical formula of Melle-Hern\'andez
\cite[\S3.4, Remark~2]{GLM97} gives the same Milnor number:
$\operatorname{Sing}V(g_e^r)=C$ is disjoint from
$\operatorname{Sing}H$, and $I_p(V(g_e^r),H)=rq_p$.
Taking $d=A$, $k=1$, and $\sum_pq_p=eD$ gives
\eqref{eq:global-mu}.

Let $Z_{\mathrm{tr}},\Delta_{\mathrm{tr}},\mu_{\mathrm{tr}}$
denote the specialization in which all $n=eD$ intersections
are transverse. The formulas for
$\zeta_{\mathrm{tr}}^{-1}=(1-t)\Delta_{\mathrm{tr}}$ and
$\mu_{\mathrm{tr}}$ agree with~\cite[\S3.5]{GLM97}, taking
their $d=A$, $k=1$, and a single smooth component of degree
$e$ and multiplicity $r$. Substituting $b=n-S$ in the
preceding theorem gives the compact form of the contact contribution:
\begin{align}
 Z_f-Z_{\mathrm{tr}}
   &=\frac{s}{(s+1)Q(s)}
       \bigl[(r-1)U_f(0)-r(Ds+4)U_f(Bs+\beta)\bigr],
       \label{eq:compressed-z}\\
 \frac{\Delta_f(t)}{\Delta_{\mathrm{tr}}(t)}
   &=\left(\frac{1-t^D}{1-t^{rD}}\right)^S
       \Gamma_f(t^B),\qquad
 \mu_f-\mu_{\mathrm{tr}}=S.\label{eq:compressed-delta}
\end{align}
Thus a contact changes both the contribution along the curve
and the local term at the intersection point. The two terms
in \eqref{eq:compressed-z} record these changes separately.
The function $U_f$ has positive simple-pole weights; for
$x>0$ and $k\ge0$,
\begin{equation}\label{eq:complete-monotonicity}
 (-1)^kU_f^{(k)}(x)
   =k!\sum_{q\ge2}
       \frac{n_q(q-1)q^k}{(1+qx)^{k+1}}\ge0.
\end{equation}

\subsection{Characters of Milnor monodromy}

To read individual eigenvalues, take a common cyclic group
$C_M$, where
$M=\operatorname{lcm}(A,rD,\{qB:n_q>0,\ q\ge2\})$.
For $N\mid M$, let $\Reg_N$ be the regular representation
of $C_N$ pulled back to $C_M$. For $qN\mid M$, set
\[
 W_{q,N}=\ker(\Reg_{qN}\longrightarrow\Reg_N),
\]
where the map sends the basis vector indexed by $j$ to the
one indexed by $j\bmod N$. A character of order $m$ has
multiplicity $\one_{m\mid qN}-\one_{m\mid N}$ in $W_{q,N}$.

\begin{proposition}\label{prop:characters}
The semisimplification of Milnor monodromy has character
\begin{equation}\label{eq:positive-character}
 [H^2(F_f,\CC)_{\mathrm{ss}}]
   =\kappa[\Reg_A]-[\one]+(b-c)[W_{r,D}]
                 +\sum_{q\ge2}n_q[W_{q,B}].
\end{equation}
For every $m\ge1$ its eigenvalue multiplicity is
\begin{equation}\label{eq:multiplicities}
 \begin{aligned}
 M_m(f)={}&\kappa\one_{m\mid A}-\one_{m=1}
   +(b-c)(\one_{m\mid rD}-\one_{m\mid D})\\
 &+\sum_{q\ge2}n_q(\one_{m\mid qB}-\one_{m\mid B}).
 \end{aligned}
\end{equation}
In particular,
\begin{equation}\label{eq:m1-mr}
 M_1(f)=2-c=2g(C),\qquad M_r(f)=b+3-2c.
\end{equation}
\end{proposition}

\begin{proof}
All roots of \eqref{eq:global-delta} are $M$th roots of unity,
so the semisimple action factors through $C_M$. Taking
characteristic polynomials gives \eqref{eq:positive-character},
and then \eqref{eq:multiplicities}. This is a positive character:
\[
 \kappa\Reg_A-\one
    =(\kappa-1)\Reg_A\oplus\Aug_A,
\]
where $\Aug_A$ is the augmentation kernel. The remaining
coefficients are nonnegative, since $b-c>0$.

For $m=1$ each $W$ term has multiplicity zero and the first
two terms give $\kappa-1=2-c$. For $m=r$, use
$r\mid A$, $r\nmid D$, and $r\mid B$, the last identity following
from $B=r(e(r-1)+1)$. This gives
$M_r=\kappa+b-c=b+3-2c$.
\end{proof}

\section{Reconstructing the contact profile}\label{sec:reconstruction}

\subsection{Recovery from monodromy coefficients}

The first step recovers the parameters before separating any
contact factors. Put
\[
 P_f(t)=(1-t)\Delta_f(t).
\]
Equation \eqref{eq:global-delta} reads
\begin{equation}\label{eq:p-factorization}
 P_f(t)=(1-t^A)^\kappa
       (1+t^D+\cdots+t^{(r-1)D})^{b-c}\Gamma_f(t^B).
\end{equation}

\begin{proposition}[Coefficient recovery]\label{prop:coefficient-recovery}
The coefficients of $\Delta_f$ determine
\begin{equation}\label{eq:parameter-recovery}
 \begin{gathered}
 A=\ord_{t=0}(P_f-1),\qquad
 \kappa=-[t^A]P_f,\qquad
 e=\frac{3+\sqrt{4\kappa-3}}2,\qquad r=\frac Ae,\\
 b=c+[t^{A+1}]P_f,\qquad
 \theta:=\sum_{q\ge2}n_q=[t^B]P_f,\qquad n_1=b-\theta.
 \end{gathered}
\end{equation}
\end{proposition}

\begin{proof}
Since $A<D<B$, the first nonconstant term in
\eqref{eq:p-factorization} is $-\kappa t^A$.
The identity $\kappa=e^2-3e+3$ gives
$4\kappa-3=(2e-3)^2$; $e\ge3$ selects the displayed root.
The coefficient of $t^D$ comes only from the second factor,
as $2A>D$ and $B>D$, and equals $b-c$.

To read the coefficient of $t^B$, consider the first two factors.
Every degree in their product has the form $iA+jD$, with
$i,j\ge0$. If it equaled $B=(r-1)A+r$, then $j\le r-1$,
whereas reduction modulo $A$ would give $j\equiv r\bmod A$.
This is impossible because $r<A$. Thus only the linear term
of $\Gamma_f(t^B)$ contributes at degree $B$, giving
$[t^B]P_f=\Gamma_f'(0)=\theta$.
\end{proof}

With $e,r,b$ known, extract $\Gamma_f$ from
\begin{equation}\label{eq:gamma-extraction}
 \frac{P_f(t)}{(1-t^A)^\kappa}
       \left(\frac{1-t^D}{1-t^{rD}}\right)^{b-c}
       =\Gamma_f(t^B).
\end{equation}
The substitution $u\mapsto t^B$ is injective, so this determines
$\Gamma_f(u)$ uniquely. Its degree is $S=n-b$.

\begin{proposition}[Finite inversion]\label{prop:finite-inversion}
Write $\Gamma_f(u)=\sum_{j=0}^S\gamma_j u^j$, with
$\gamma_0=1$, and extend $\gamma_j=0$ for $j>S$. Define
\begin{equation}\label{eq:log-recurrence}
 L_m=m\gamma_m-\sum_{j=1}^{m-1}\gamma_jL_{m-j}.
\end{equation}
Then the entire profile is recovered by
\begin{equation}\label{eq:mobius-recovery}
 n_q=-\frac1q\sum_{d\mid q}\mob(q/d)L_d\quad(q\ge2),
 \qquad n_1=b-\gamma_1.
\end{equation}
It suffices to compute $L_m$ for $m\le S+1$.
\end{proposition}

\begin{proof}
Recurrence \eqref{eq:log-recurrence} is coefficient comparison
in
$\Gamma_f(u)\sum_{m\ge1}L_m u^m=u\Gamma_f'(u)$.
On the other hand, the product in \eqref{eq:contact-transforms}
gives
\[
 \frac{u\Gamma_f'(u)}{\Gamma_f(u)}
   =\sum_{q\ge2}n_q
       \left(\frac{u}{1-u}-\frac{q u^q}{1-u^q}\right).
\]
Consequently
$L_m=\theta-\sum_{q\mid m,\ q\ge2}qn_q$.
For $q>1$, M\"obius inversion cancels the constant $\theta$
and gives \eqref{eq:mobius-recovery}. Finally, an occurring
$q\ge2$ satisfies $q-1\le S$. If $S=0$, then $\Gamma_f=1$
and all intersections are transverse.
\end{proof}

Appendix~\ref{app:example} carries out this recovery for an
integer-coefficient member of the family.

\paragraph{Lefschetz numbers.}
The multiplicity of the surface germ is also determined by
its Lefschetz numbers. Let
$\Lambda_m=1+\Tr(T_f^m\mid H^2(F_f,\CC))$. The formal identity
\[
 -t\frac{P_f'(t)}{P_f(t)}
       =\sum_{m\ge1}\Lambda_m t^m
\]
and the initial term $-\kappa t^A$ give
$\Lambda_m=0$ for $1\le m<A$ and $\Lambda_A=A\kappa$.
Thus $A$ is the first index with a nonzero Lefschetz number.

\subsection{Recovery from poles and residues}

To reconstruct from $Z_f$, we locate its poles and compute
their residues. The residue at the largest pole distinguishes
the cubic case from degrees $e\ge4$.

\begin{proposition}\label{prop:poles}
All actual poles of $Z_f$ are simple. Its largest pole
$\sigma_*$ and corresponding residue $R_*$ are
\begin{equation}\label{eq:largest-pole}
 (\sigma_*,R_*)=
 \begin{cases}
  (-3/A,-3(e-1)/A),&e\ge4,\\
  (-1/r,4/r),&e=3.
 \end{cases}
\end{equation}
Every occurring $q\ge2$ gives a distinct actual pole
$\sigma_q=-(q\beta+1)/(qB)$, with
\begin{equation}\label{eq:contact-residue}
 \Res_{s=\sigma_q}Z_f
  =-\frac{n_q(q-1)r\sigma_q(D\sigma_q+4)}
          {qB(\sigma_q+1)Q(\sigma_q)}\ne0.
\end{equation}
If $e\ge4$, then $\sigma_Q=-(3r+1)/(rD)$ is also an actual
pole with positive residue.
\end{proposition}

\begin{proof}
Formula \eqref{eq:global-z} gives the candidate list
\begin{equation}\label{eq:candidates}
 -1,\quad -3/A,\quad -(3r+1)/(rD),\quad
 \sigma_q\quad(q\ge2,\ n_q>0).
\end{equation}
\emph{Degrees $e\ge4$.}
The strict inequalities
\begin{equation}\label{eq:pole-order}
 0<\frac3A<\frac{3r+1}{rD}<\frac{\beta}{B}
      <\frac{q\beta+1}{qB}<1
\end{equation}
show that the candidate roots are distinct, apart from repeated
occurrences of the same contact order. The numerators of the
first two differences are $r(e-3)$, and
$B-\beta=(r-1)(D-3)$ gives the last inequality.
Thus every actual pole is simple. At $s=-3/A$,
$Q(s)=(e-3)/e$, so the residue is
\[
 \frac{\kappa+c/Q(-3/A)}A=-\frac{3(e-1)}A.
\]
This proves the first case of \eqref{eq:largest-pole}.

The residue at a contact root is \eqref{eq:contact-residue}.
All its factors are nonzero except possibly $D\sigma_q+4$.
If that factor vanished, then
\[
 D=q\bigl[(r-2)D+4\bigr].
\]
For $r\ge3$ the right hand side is larger than $D$.
For $r=2$ this would say that the odd integer $D=2e+1$
equals $4q$. Both are impossible.

At the root $\sigma_Q$ of $Q$, set
\[
 \delta=\frac{e-3}{D},\qquad
 \eta=\frac{(3r+1)(r-1)}{rD[rD-(3r+1)]}.
\]
At $\sigma_Q$ we have $Bs+\beta=\delta$ and $rJ=r-1$.
The transverse residue, obtained from \eqref{eq:global-z}
with $b=n$, is $\eta e(2D+e-3)$. By
\eqref{eq:compressed-z}, the contact correction has residue
$-\eta[S-U_f(\delta)]$. Consequently
\begin{equation}\label{eq:old-positive-residue}
 \Res_{s=\sigma_Q}Z_f
    =\eta\bigl[e(D+e-3)+b+U_f(\delta)\bigr]>0.
\end{equation}

\emph{The cubic case.}
Now $c=0$, $B=r\beta$, and
\[
 a=3(rs+1),\qquad Q=D(rs+1),\qquad D=3r+1.
\]
The sole term in \eqref{eq:global-z} with denominator $aQ$
has coefficient zero. More explicitly, the formula simplifies
to
\begin{equation}\label{eq:cubic-formula}
 Z_f(s)=\frac4{rs+1}-\frac3{s+1}
       +\sum_{q\ge2}n_q\frac{C_qs}{(s+1)K_q(s)},
 \qquad
 C_q=\frac{(q-1)[q(r-1)\beta-D]}D>0.
\end{equation}
To obtain it, subtract the transverse formula and use
$(r-1)K_q-rJ=[q(r-1)\beta-D](rs+1)$.
For positivity, the smallest possible bracket is bounded below
by $6r^2-13r+3$, which equals $1$ at $r=2$ and increases for
$r\ge2$.

The contact roots are now
$\sigma_q=-1/r-1/(qB)\in(-1,-1/r)$ and $K_q(-1/r)=1$.
The common root $-1/r$ has residue $4/r$, while every contact
residue has the simpler expression
\begin{equation}\label{eq:cubic-contact-residue}
 \Res_{s=\sigma_q}Z_f
       =\frac{n_qC_q\sigma_q}{qB(\sigma_q+1)}<0.
\end{equation}
Finally, using $\sum qn_q=3D$, the residue at $-1$ is
\begin{equation}\label{eq:cubic-minus-one}
 \Res_{s=-1}Z_f
   =-\frac1D\left[
       b+3r\sum_{q\ge2}
          \frac{n_q(q-1)}{q(r-1)\beta-1}\right]<0.
\end{equation}
This proves all assertions at the cubic boundary as well.
\end{proof}

The remaining candidate $-1$ has residue
\begin{equation}\label{eq:minus-one-residue}
 \Res_{s=-1}Z_f
 =\frac{\displaystyle c+b(r-1)+r(D-4)
       \sum_{q\ge2}\frac{n_q(q-1)}
                  {q(r-1)(D-3)-1}}
        {3r+1-rD}.
\end{equation}
Together with Proposition~\ref{prop:poles}, this determines the
pole set: $-1$ is present exactly when this residue is nonzero.

\begin{proposition}\label{prop:z-recovery}
The rational function $Z_f$ determines $e,r$ and the contact
profile.
\end{proposition}

\begin{proof}
Take the largest actual pole $\sigma_*$ and its residue $R_*$.
Proposition~\ref{prop:poles} distinguishes the two parameter
ranges by the sign of $R_*$:
\begin{equation}\label{eq:z-parameter-recovery}
 \begin{cases}
  R_*>0:\quad e=3,\qquad r=-1/\sigma_*;\\
  R_*<0:\quad e=1+R_*/\sigma_*,
           \qquad r=-3/(e\sigma_*).
 \end{cases}
\end{equation}
Now all the numerical constants in
\eqref{eq:linear-functions} are known. Remove the positions
$-1,-3/A,-(3r+1)/(rD)$ from the pole set. Its remaining
elements are exactly the occurring contact poles.
At such a pole $\sigma_q$ with residue $R_q$, use
\begin{equation}\label{eq:z-contact-recovery}
 q=-\frac1{B\sigma_q+\beta},\qquad
 n_q=-\frac{qB(\sigma_q+1)Q(\sigma_q)}
         {(q-1)r\sigma_q(D\sigma_q+4)}\,R_q.
\end{equation}
The denominators are nonzero by
Proposition~\ref{prop:poles}. Finally
$n_1=eD-\sum_{q\ge2}qn_q$.
\end{proof}

\begin{proof}[Proof of Theorem~\ref{thm:main}]
Propositions~\ref{prop:coefficient-recovery}
and~\ref{prop:finite-inversion} recover all the data from
$\Delta_f$, and Proposition~\ref{prop:z-recovery} recovers them
from $Z_f$. Conversely, the contact formulas determine both
invariants from those data. This proves the three equivalent
conditions and the stated recovery procedures. Over $\CC$,
the characteristic polynomial determines a linear
automorphism up to semisimplification, which proves the
classification statement.
\end{proof}

\subsection{Monodromy and analytic moduli}

\begin{corollary}\label{cor:monodromy}
Every actual pole $s_0$ of $Z_f$ gives the eigenvalue
$\exp(2\pi i s_0)$ on the cohomology of $F_f$.
\end{corollary}

\begin{proof}
Use the positive character \eqref{eq:multiplicities}.
For $s_0=-3/A$, the eigenvalue order is
$m=A/\gcd(A,3)>1$, detected by $\Reg_A$.
For $s_0=-(3r+1)/(rD)$, it is
$m=rD/\gcd(rD,3r+1)>1$. If $m\mid D$, then
$rD\mid(3r+1)D$, which would force $r\mid3r+1$.
Thus $W_{r,D}$ detects this order.
For a contact pole, the order is
$m=qB/\gcd(q\beta+1,qB)$. If $m\mid B$, then
$qB\mid(q\beta+1)B$, contradicting $q\nmid q\beta+1$.
It is therefore detected by $W_{q,B}$.
The order at $-1$ is $1$, with multiplicity $M_1=2-c>0$.
These cases also cover the common old root when $e=3$.
All primitive roots of each detected order occur with the
same multiplicity by \eqref{eq:multiplicities}.
\end{proof}

\begin{example}[Different curves with the same invariants]
\label{ex:cubic-moduli}
Take two smooth plane cubics with different $j$-invariants,
fix the same $r\ge2$, and choose a degree $3r+1$ curve transverse
to each. Both profiles consist of $3(3r+1)$ ones, so the
corresponding $Z_f$ and $\Delta_f$ agree.
The resulting germs are not analytically contact equivalent.
Such an equivalence would identify their lowest homogeneous
terms up to a linear change of coordinates and a nonzero
scalar. Their tangent-cone supports would then be projectively
equivalent cubics, contrary to their distinct $j$-invariants.
\end{example}

\section{Spectral balance and contact collisions}\label{sec:collisions}

\subsection{The balance identity}

\begin{proof}[Proof of Theorem~\ref{thm:balance}]
Equation~\eqref{eq:m1-mr} gives $M_r-M_1=b+1-c$.
Adding this to the second expression in \eqref{eq:global-mu}
proves \eqref{eq:balance-intro}.
Fix a primitive $r$th root $\lambda$. An equivariant injection
is injective on the generalized $\lambda$-eigenspaces, so both
the total dimension $\mu$ and the dimension $M_r$ of this
individual generalized eigenspace can only increase.
At fixed $e,r$, the balance identity
and the fixed value $M_1=2-c$ force equality in both dimensions.
The injection is therefore an isomorphism. The argument also
applies to semisimplifications. If the contact profiles differ,
Theorem~\ref{thm:main} gives distinct characteristic polynomials,
so neither direction admits such an injection.
\end{proof}

\begin{example}[Equal Milnor numbers and incomparable monodromy]
\label{ex:equal-mu}
The profiles $\{2,4,1^{30}\}$ and $\{3,3,1^{30}\}$ both occur
with $e=4,r=2$, even on the same smooth quartic $C$.
To realize them, choose distinct $p_1,p_2\in C$ and prescribe
$\mathfrak d=2p_1+4p_2$ or $3p_1+3p_2$.
The line bundle
$\mathcal L=\mathcal O_C(9)(-\mathfrak d)$ has degree $30$
on a curve of genus $3$. Riemann--Roch gives
$h^0(\mathcal L)=28$. Vanishing at either $p_i$ is one
independent condition, and a double zero at any other point
$x$ imposes two independent conditions, since
$H^1(\mathcal L(-2x))=0$.
The double-zero incidence has dimension at most $27$ as
$x$ varies in $C$. A general section thus has $30$ simple
zeros away from $p_1,p_2$, giving the desired divisor on $C$.

The restriction of degree-nine homogeneous polynomials to
$C$ is surjective, because $H^1(\PP^2,\mathcal O(5))=0$.
After lifting the section, replace its lift $h_9$ by
$h_9+g_4k_5$. Degree-five polynomials $k_5$ can adjust the
two normal derivatives independently. Choose the resulting
derivatives nonzero at both prescribed contacts. At the other
intersections, $h_9|_C$ has simple zeros. Thus $H$ is smooth
at every intersection and both profiles occur in
\eqref{eq:family}.

Their contact polynomials are
\[
 \Gamma_{2,4}(u)=1+2u+2u^2+2u^3+u^4,\qquad
 \Gamma_{3,3}(u)=1+2u+3u^2+2u^3+u^4.
\]
Both have $b=32$, $M_1=6$, and the remaining data are
\[
 \begin{array}{c|rrrr}
       &\mu&M_2&M_{20}&M_{30}\\ \hline
 f_{2,4}&419&43&2&0\\
 f_{3,3}&419&43&0&2
 \end{array}
 \qquad
 \frac{\Delta_{2,4}(t)}{\Delta_{3,3}(t)}
       =\frac{(1-t^{20})(1-t^{40})}{(1-t^{30})^2}.
\]
The two representations have equal dimension and admit no
equivariant injection in either direction.
\end{example}

\subsection{Contact collisions}

\begin{proposition}[Collision formulas]\label{prop:collision}
Consider two members with the same $e,r$ whose profiles differ
by replacing $q$ transverse intersections by one intersection
of multiplicity $q\ge2$. Denote these members by
$f_{\mathrm{split}}$ and $f_{\mathrm{tan}}$, respectively.
Then
\begin{align}
 \frac{\Delta_{\mathrm{tan}}(t)}
      {\Delta_{\mathrm{split}}(t)}
   &=\frac{1-t^{qB}}{1-t^B}
        \left(\frac{1-t^D}{1-t^{rD}}\right)^{q-1},
        \label{eq:collision-delta}\\
 [H^2_{\mathrm{tan,ss}}]-[H^2_{\mathrm{split,ss}}]
   &=[W_{q,B}]-(q-1)[W_{r,D}],\label{eq:collision-character}\\
 \mu_{\mathrm{tan}}-\mu_{\mathrm{split}}&=q-1,\qquad
 M_{r,\mathrm{tan}}-M_{r,\mathrm{split}}=-(q-1).
        \label{eq:collision-dimensions}
\end{align}
The topological functions satisfy
\begin{equation}\label{eq:collision-z}
 Z_{\mathrm{tan}}-Z_{\mathrm{split}}
   =\frac{(q-1)(r-1)s}{dQ}
        -\frac{(q-1)rsJ}{dQK_q}.
\end{equation}
\end{proposition}

\begin{proof}
The collision decreases $b$ by $q-1$, adds one factor
$(1-t^{qB})/(1-t^B)$, and adds $(q-1)/K_q$ to
$U_f(Bs+\beta)$. Substitution into the global formulas proves
\eqref{eq:collision-delta}, \eqref{eq:collision-character},
and \eqref{eq:collision-z}. Taking dimensions and using
$B-(r-1)D=1$ gives \eqref{eq:collision-dimensions}.
\end{proof}

The character difference records the new block $W_{q,B}$ and
the decrease of the $W_{r,D}$ contribution. The net gain in
dimension is $q-1$, while $M_r$ drops by $q-1$.

\begin{proposition}[Splitting a single contact]\label{prop:splitting}
Suppose that $H_0$ has one contact of order $q\ge2$ with $C$,
at $p$, and all its other intersections with $C$ are transverse.
Choose a linear form $L$ with $L(p)\ne0$ and set
\begin{equation}\label{eq:splitting-family}
 h_{D,\tau}=h_{D,0}+\tau L^D,\qquad
 f_\tau=g_e^r+h_{D,\tau}.
\end{equation}
For every sufficiently small $\tau\ne0$, the point $p$ splits
into $q$ transverse intersections, all other intersections
remain transverse, and $f_\tau$ satisfies the assumptions
of \eqref{eq:family}.
\end{proposition}

\begin{proof}
Near $p$ use a parameter $w$ on $C$, with $L$ nonvanishing.
The function $\phi=-h_{D,0}|_C/L^D|_C$ has a zero of order
$q$ at $w=0$, and $\phi'$ is $w^{q-1}$ times a unit.
Shrink the disc so that $0$ is its only critical point.
For every sufficiently small nonzero $\tau$, the equation
$\phi(w)=\tau$ has $q$ roots by Rouch\'e's theorem, and all
are simple because their derivatives do not vanish.

The remaining simple intersections persist by the implicit
function theorem. Compactness of $C$ prevents additional
zeros outside disjoint neighborhoods of the old intersections.
At $p$ the normal derivative of $h_{D,0}$ along $C$ is nonzero,
and remains so near the new roots. At all other roots the
tangential derivative is nonzero. Thus $H_\tau$ is smooth at
every intersection, with all the claimed intersection
multiplicities. Lemma~\ref{lem:isolated} applies to each $f_\tau$.
\end{proof}

\subsection{Extremal contact profiles}

\begin{proposition}[Extremal Milnor numbers]\label{prop:milnor-bounds}
At fixed $e,r$,
\begin{equation}\label{eq:milnor-bounds}
 \mu_{\mathrm{tr}}\le\mu_f\le\mu_{\mathrm{tr}}+eD-1,
 \qquad
 \mu_{\mathrm{tr}}=A\kappa+(r-1)D(eD-c)-1.
\end{equation}
Both bounds are attained for every $e\ge3,r\ge2$.
Equality at the lower bound means that all intersections are
transverse; equality at the upper bound means that there is
a single intersection of multiplicity $eD$.
\end{proposition}

\begin{proof}
The classical Milnor formula \eqref{eq:global-mu} gives
$\mu_f-\mu_{\mathrm{tr}}=n-b$, and $1\le b\le n$.
For the lower bound take $h_D$ to be a product of $D$ general
linear forms. Their intersections with $C$ are transverse and
pairwise distinct on $C$, and $H$ is smooth at each of them.

For the upper bound choose
\[
 g_e=x^e+y^e+z^e,\qquad
 \alpha^e=-1,\qquad p=[1:\alpha:0],\qquad
 \ell=x+\alpha^{e-1}y.
\]
On $\ell=0$ one has $g_e=z^e$, so
$\operatorname{div}_C(\ell)=e\,p$. Set
\begin{equation}\label{eq:maximal-contact-example}
 h_D=\ell^D+g_e x^{D-e}.
\end{equation}
Its restriction to $C$ has divisor $eD\,p$.
In the affine chart $x=1$,
$dh_D(p)=dg_e(p)\ne0$, so $H$ is smooth at its unique
intersection with $C$. This realizes the upper bound.
\end{proof}

The maximal-contact member \eqref{eq:maximal-contact-example}
has exactly one intersection, so it supplies such a family
for every $e,r$, with $q=eD$.

\section{Twisted topological zeta functions}\label{sec:twisted}

\subsection{The divisibility-filtered formula}

Fix an integer $\ell>1$. On a log resolution define
\begin{equation}\label{eq:twisted-definition}
 Z_f^{(\ell)}(s)=
   \sum_{\substack{\varnothing\ne I\\
                  \ell\mid N_i\ \text{for every }i\in I}}
       \chi(E_I^\circ\cap h^{-1}(0))
            \prod_{i\in I}(N_i s+\nu_i)^{-1}.
\end{equation}
This is the twisted topological normalization of
\cite[Section~1.3, equation~(1.6)]{AGGL26}.
Write $\delta(N)=\one_{\ell\mid N}$ and $\eta_q=\gcd(D,q)$.

The auxiliary coordinate boundaries $\mathsf V,\mathsf W$
have $N=0$, so the condition $\ell\mid0$ preserves the Euler
refinement used in the local calculations.

\begin{theorem}\label{thm:twisted-formula}
Set
\begin{equation}\label{eq:twisted-e}
 E_{q,\ell}(s)=q\delta(r)+r\delta(\eta_q)
       -\delta(r\eta_q)
       +Q(s)\bigl[\delta(qB)-q\delta(B)\bigr].
\end{equation}
The complete twisted function is
\begin{equation}\label{eq:global-twisted}
 Z_f^{(\ell)}
 =\frac{\kappa\delta(A)}a+\frac{c\delta(r)}{aQ}
   +\frac{(c-b)[r\delta(D)-\delta(rD)]}{Q}
   +\sum_{q\ge2}\frac{n_q E_{q,\ell}}{QK_q}.
\end{equation}
\end{theorem}

\begin{proof}
The filter in \eqref{eq:twisted-definition} is compatible with
the normal crossings change in \eqref{eq:snc-identity}.
If the old multiplicities at a blowup center are all divisible
by $\ell$, the new one, their sum, is also divisible, and the
same identity applies. If exactly one is not divisible, the
new multiplicity is not divisible and all contributions vanish.
If at least two are not divisible, every projective stratum
of nonzero Euler characteristic retains a nondivisible old
component, so again every contribution vanishes. Thus the
filtered local functions also assemble by Euler Fubini.

On the regular toric resolution, the strict transform has
$N=1$ and is excluded. The auxiliary boundaries just described
are retained for the orbit stratification. Filtering the orbit
calculation by the saturated-lattice gcds in
Appendix~\ref{app:filter} gives
\begin{equation}\label{eq:local-twisted}
 \begin{aligned}
 Z_{u^A}^{(\ell)}&=\frac{\delta(A)}a,\\
 Z_G^{(\ell)}&=\frac{\delta(r)}{aQ}
                  +\frac{r\delta(D)-\delta(rD)}Q,\\
 Z_q^{(\ell)}&=\frac{\delta(r)}{aQ}
                  +\frac{E_{q,\ell}}{QK_q}.
 \end{aligned}
\end{equation}
When $q=1$, $\eta_q=1$ and $E_{1,\ell}=0$.
Apply \eqref{eq:assembly-z} with these local functions;
the coefficient of $\delta(r)/(aQ)$ is $c$, giving
\eqref{eq:global-twisted}.
\end{proof}

\subsection{Holomorphy}

\begin{theorem}\label{thm:holomorphy}
If no monodromy eigenvalue at the origin has order divisible
by $\ell>1$, then $Z_f^{(\ell)}$ is identically zero.
\end{theorem}

\begin{proof}
By \eqref{eq:positive-character}, every eigenvalue order
divides one of
\[
 A,\qquad rD,\qquad \{qB:q\ge2,\ n_q>0\}.
\]
Conversely, each of these integers itself occurs as an
eigenvalue order: primitive $A$th, $rD$th, and $qB$th
characters have positive multiplicity in the corresponding
positive representation blocks. The hypothesis is therefore
equivalent to
\[
 \ell\nmid A,\qquad \ell\nmid rD,\qquad
 \ell\nmid qB\quad(q\ge2,\ n_q>0).
\]
It gives $\delta(A)=\delta(rD)=0$, and the divisibilities
$r\mid A$, $D\mid rD$, $\eta_q\mid D$,
$r\eta_q\mid rD$ give
$\delta(r)=\delta(D)=\delta(\eta_q)
=\delta(r\eta_q)=0$.
For each occurring $q\ge2$ it also gives
$\delta(qB)=\delta(B)=0$.
Every term in \eqref{eq:global-twisted} vanishes.
\end{proof}

At smooth points the monodromy has eigenvalue $1$, so the
same criterion applies to eigenvalues in a neighborhood
of the origin.

\begin{example}[The order condition]\label{ex:twisted-order}
For $e=4,r=2$ with all intersections transverse, one has
$c=-4,b=36$ and
\[
 M_3(f)=0,\qquad Z_f^{(3)}(s)=-\frac{40}{18s+7}.
\]
A primitive $18$th root is nevertheless an eigenvalue, and
$3\mid18$.
\end{example}

\begin{corollary}\label{cor:twisted-recovery}
Within the family \eqref{eq:family}, the single polynomial
$\Delta_f$ determines every twisted function $Z_f^{(\ell)}$.
\end{corollary}

\begin{proof}
Recover $e,r,\{q_p\}$ by Theorem~\ref{thm:main}, and use
\eqref{eq:global-twisted}.
\end{proof}

\appendix
\section{Weighted cone calculations}\label{app:cones}

\subsection{Orbit contributions and subdivision}\label{app:subdivision}

We derive the weighted formula of Proposition~\ref{prop:local}.
For a nonnegative vector $\rho$ set
\[
 N(\rho)=\min\{\langle X,\rho\rangle,\langle Y,\rho\rangle,
                      \langle Z,\rho\rangle\},\qquad
 \nu(\rho)=3\rho_u+\rho_v+\rho_w.
\]
For a regular toric cone with primitive rays $\rho_i$,
the monomial substitution gives the actual numerical data
$N(\rho_i),\nu(\rho_i)$. The reduced strict transform
$\mathcal S$ has data $(1,1)$, and the auxiliary boundaries
$\mathsf V,\mathsf W$ have data $(0,1)$.

Let $\sigma$ be a regular cone of dimension $k$, and let $\tau$
be the face selected by its relative interior. Its orbit maps
to the origin exactly when $\relint(\sigma)$ contains a vector
with all coordinates positive. Thus all relevant selected
faces are compact. If $k=3$, the selected face is a vertex
and the orbit is a point disjoint from the strict transform.
Its contribution is $\prod_i(N(\rho_i)s+\nu(\rho_i))^{-1}$.

For $k<3$, the orbit is $(\CC^*)^{3-k}$. If
$\dim\tau<3-k$, the strict-transform equation either has no
zero or has an additional torus factor. Both the orbit and its
strict-transform zero set have Euler characteristic zero.
To see the factor directly, the edge difference vector is
primitive and the triangle difference lattice is saturated,
as their normalized volumes are one. They can therefore be
completed to a basis in the orbit character lattice.

If $\dim\tau=3-k\ge1$, the zero set is one point for an edge,
and $1+\xi+\eta=0$ in $(\CC^*)^2$ for the triangle.
Writing $d_\tau=\dim\tau$, its Euler characteristic is
$(-1)^{d_\tau-1}$. The complement and strict-transform terms
together contribute
\[
 \left[-\chi(\mathcal S\cap O_\sigma)
      +\frac{\chi(\mathcal S\cap O_\sigma)}{s+1}\right]
      \prod_i(N(\rho_i)s+\nu(\rho_i))^{-1}
 =\frac{s}{s+1}(-1)^{d_\tau}
      \prod_i(N(\rho_i)s+\nu(\rho_i))^{-1}.
\]

These terms can be merged over nonregular cones while keeping
their lattice multiplicities. For a simplicial cone $\Delta$,
write
\[
 J_\Delta(s)=
   \frac{\mult(\Delta)}
        {\prod_{\rho_i\text{ ray of }\Delta}
                    (N(\rho_i)s+\nu(\rho_i))}.
\]
The multiplicity is the index of its rays in the saturated
lattice of its span. On a Newton cone $N$ is linear, and,
for real $s>0$,
\[
 J_\Delta(s)=\int_\Delta e^{-sN(x)-\nu(x)}\,dx,
\]
with the lattice-normalized measure on that span. The integral
is additive under subdivisions of the same dimension.
Equality of the resulting rational functions then follows
by continuation. In particular, $J_\tau$, the sum over a
triangulation of the dual cone of $\tau$, is independent of
that triangulation. The preceding orbit calculation yields
\begin{equation}\label{eq:weighted-newton}
 Z_q=J_X+J_Y+J_Z+
       \frac{s}{s+1}(-J_{XY}-J_{XZ}-J_{YZ}+J_{XYZ}).
\end{equation}
This is the Newton face formula specialized to the pulled-back
volume form; compare the standard-volume version
in~\cite[Th\'eor\`eme~5.3]{DL92}.

\subsection{The cone table}

Using the normals in the proof of Proposition~\ref{prop:local},
triangulate the vertex cones as follows:
\begin{center}
\begin{tabular}{lll}
\toprule
Vertex & Simplicial cones & Multiplicities\\
\midrule
$X$ & $(\mathsf U,\mathsf W,\mathsf T)$,
       $(\mathsf U,\mathsf R,\mathsf T)$ & $q,\ 1$\\
$Y$ & $(\mathsf W,\mathsf S,\mathsf T)$ & $q^2(r-1)$\\
$Z$ & $(\mathsf V,\mathsf R,\mathsf T)$,
       $(\mathsf V,\mathsf T,\mathsf S)$ & $r,\ q(r-1)$\\
\bottomrule
\end{tabular}
\end{center}
Absolute determinants give the last column. Hence
\[
 J_X=\frac q{aK_q}+\frac1{aQK_q},\qquad
 J_Y=\frac{q^2(r-1)}{p_qK_q},\qquad
 J_Z=\frac r{QK_q}+\frac{q(r-1)}{p_qK_q}.
\]
The two-dimensional lattice indices and the primitive ray
$\mathsf T$ give
\[
 J_{XY}=\frac q{K_q},\qquad
 J_{XZ}=\frac1{QK_q},\qquad
 J_{YZ}=\frac{q(r-1)}{p_qK_q},\qquad
 J_{XYZ}=\frac1{K_q}.
\]
In \eqref{eq:weighted-newton}, all terms containing $p_q$ combine
as
\[
 \frac{q(r-1)}{p_qK_q}
       \left(q+1-\frac{s}{s+1}\right)
       =\frac{q(r-1)}{(s+1)K_q}.
\]
Substituting the remaining terms gives
\eqref{eq:local-raw}. This identifies the cancellation of the
apparent denominator $p_q$ before the global assembly.

\subsection{Divisibility on saturated lattices}\label{app:filter}

For a compact face $\tau$, let $C_\tau$ be its dual cone and
$L_\tau=\operatorname{span}_{\RR}(C_\tau)\cap\ZZ^3$.
Choose an exponent $v_\tau\in\tau$ and put
\[
 m_\tau=\gcd\{\langle v_\tau,\rho\rangle:\rho\in L_\tau\}.
\]
The definition is independent of the exponent chosen in the
face, since all exponent differences annihilate $L_\tau$.
It uses the linear extension of $N|_{C_\tau}$ to this lattice.
The primitive rays of a relatively full-dimensional regular
subcone of $C_\tau$ form a basis of $L_\tau$. Thus all their
multiplicities are divisible by $\ell$ exactly when
$\ell\mid m_\tau$, so the divisibility indicator is constant
on these subcones. By the lattice-normalized subdivision
additivity in Appendix~\ref{app:subdivision}, their filtered
cone sum is $\delta(m_\tau)J_\tau$.

For the present triangle these gcds are
\begin{center}
\begin{tabular}{lrrrrrrr}
\toprule
$\tau$ & $X$ & $Y$ & $Z$ & $XY$ & $XZ$ & $YZ$ & $XYZ$\\
\midrule
$m_\tau$ & $r$ & $1$ & $\eta_q$ & $B$ &
            $r\eta_q$ & $\eta_q$ & $qB$\\
\bottomrule
\end{tabular}
\end{center}
The vertex rows are the gcds of their coordinates.
For $XY$, a lattice basis is $(r-1,1,0),(0,0,1)$,
whose pairings are $B,0$.
For $XZ$, the rays $\mathsf R,\mathsf T$ are a lattice basis
because their cross product is primitive. Thus
$m_{XZ}=\gcd(rD,qB)=r\gcd(D,q)$; here
$B/r=e(r-1)+1$ is coprime to $D$.
For $YZ$ use $(1,0,0),(0,q,1)$, with pairings $D,q$.
The triangle lattice is generated by $\mathsf T$, of
multiplicity $qB$.

For $\ell>1$ the strict-transform terms are excluded, so the
orbit signs are $+1$ for vertices, $-1$ for edges, and $+1$
for the triangle. Since $\delta(1)=0$, the filtered sum is
\begin{equation}\label{eq:filtered-cones}
 \begin{aligned}
 Z_q^{(\ell)}={}&\delta(r)J_X+\delta(\eta_q)J_Z
      -\delta(B)J_{XY}-\delta(r\eta_q)J_{XZ}\\
      &-\delta(\eta_q)J_{YZ}+\delta(qB)J_{XYZ}.
 \end{aligned}
\end{equation}
The two $p_q$ terms have the same indicator and cancel.
Finally
\[
 \frac q{aK_q}+\frac1{aQK_q}
       =\frac1{aQ}+\frac q{QK_q}
\]
by $K_q=q(Q-a)+1$, giving the $Z_q^{(\ell)}$ formula
in \eqref{eq:local-twisted}.
For $G$, the two vertex gcds are $r,D$ and the edge gcd
is $rD$; the same two-dimensional calculation gives
$\delta(r)/(aQ)+(r\delta(D)-\delta(rD))/Q$.

\section{An integer-coefficient example}\label{app:example}

Consider
\begin{equation}\label{eq:integer-example}
 \begin{gathered}
  g=x^4+y^4+z^4,\qquad p_7=x^7+y^7+z^7,\qquad
  p_5=x^5+y^5+z^5,\\
  h_9=z^2p_7+gp_5,\qquad f=g^2+h_9.
 \end{gathered}
\end{equation}
\paragraph{Contact profile.}
On the smooth quartic
$C=V(g)$,
$\operatorname{div}_C(h_9)=2\operatorname{div}_C(z)
                              +\operatorname{div}_C(p_7)$.

The four points with $z=0$ are $[1:t:0]$, $t^4=-1$,
and the line meets $C$ transversely.
At these points $p_7=1-t^3\ne0$ and $p_5=1-t\ne0$.
Thus $h_9|_C$ has order exactly two and
$dh_9=p_5\,dg\ne0$.

The intersection of $C$ and $V(p_7)$ has no point with a zero
coordinate. At a nontransverse intersection their gradients
are proportional, forcing $x^3=y^3=z^3$.
Set $x=1$. The conditions reduce to
$y^3=z^3=1$ and $1+y+z=0$, so the only such points are
\[
 [1:\omega:\omega^2],\qquad[1:\omega^2:\omega],
 \qquad \omega^2+\omega+1=0.
\]
Their contact orders are exactly two.
At the first point $\nabla p_7=(7/4)\nabla g$;
for the tangent vector $v=(0,1,-1)$,
\[
 v^{\mathsf t}
    \left(\operatorname{Hess}(p_7)
           -\frac74\operatorname{Hess}(g)\right)v
       =21(\omega^2+\omega)=-21\ne0.
\]
The second point is identical by symmetry. At both,
$p_5=0$ and $dh_9=z^2dp_7\ne0$.

All remaining intersections of $C$ and $V(p_7)$ are transverse.
There is no common component, since $p_7|_C$ is nonzero at
the four points with $z=0$. B\'ezout gives total intersection
number $28$; after the two double points this leaves
$24$ transverse points, all with $z\ne0$. At these points
$h_9|_C$ has a simple zero, so $H$ is smooth there as well.
Consequently
\[
 n_2=6,\qquad n_1=24,\qquad b=30,\qquad
 e=4,\quad r=2,\quad A=8,\quad D=9,\quad B=10.
\]
\paragraph{The two invariants.}
The contact formulas give
\[
 \begin{aligned}
 Z_f(s)&=\frac{-3784s^3-1922s^2+321s+189}
                {(8s+3)(s+1)(18s+7)(20s+9)},\\
 \Delta_f(t)&=\frac{(1-t^8)^7}{1-t}
                   (1+t^9)^{34}(1+t^{10})^6,\qquad
 \mu_f=421.
 \end{aligned}
\]
The contact pole $-9/20$ has residue $27/1210$.
The associated primitive $20$th roots have multiplicity six.
The balance identity reads $421+41-6=456=8^3-8^2+8$.

\paragraph{Reconstruction.}
Starting from the characteristic polynomial, the product
$P_f=(1-t)\Delta_f$ begins
\[
 P_f(t)=1-7t^8+34t^9+6t^{10}+O(t^{16}).
\]
Proposition~\ref{prop:coefficient-recovery} gives $A=8$,
$\kappa=7$, hence $e=4$, $r=2$, $c=-4$, and $b=-4+34=30$.
Removing the two determined factors leaves
\[
 \frac{P_f(t)}{(1-t^8)^7(1+t^9)^{34}}
       =(1+t^{10})^6=\Gamma_f(t^{10}).
\]
Thus $\Gamma_f(u)=(1+u)^6$. Finite inversion gives
$n_2=6$ and $n_q=0$ for $q\ge3$; finally $n_1=30-6=24$.

\section*{Acknowledgments}

This mathematical paper was prepared with the assistance of artificial intelligence.  I
first acknowledge the textual data produced by humanity as a whole, which
provides the linguistic and cultural record from which mathematical expression
can be learned.  I then thank the accumulated scholarly achievements of
mathematicians throughout history.  I also thank the researchers whose work has
advanced and refined neural networks through successive generations.  Finally,
I thank the OpenAI team for the training and development of the models used in
the preparation of this paper.

\begin{thebibliography}{99}

\bibitem{AC75}
N.~A'Campo,
\emph{La fonction z\^eta d'une monodromie},
Comment. Math. Helv. \textbf{50} (1975), 233--248.

\bibitem{AGGL26}
E.~Artal Bartolo, P.~D.~Gonz\'alez P\'erez,
M.~Gonz\'alez Villa, and E.~Le\'on Cardenal,
\emph{Denef--Loeser zeta functions of suspensions and
L\^e--Yomdin singularities},
\href{https://arxiv.org/abs/2604.24523v1}{arXiv:2604.24523v1}
(2026).

\bibitem{DL92}
J.~Denef and F.~Loeser,
\emph{Caract\'eristiques d'Euler--Poincar\'e, fonctions z\^eta
locales et modifications analytiques},
J. Amer. Math. Soc. \textbf{5} (1992), no.~4, 705--720.

\bibitem{GLM97}
S.~M.~Gusein-Zade, I.~Luengo, and A.~Melle-Hern\'andez,
\emph{Partial resolutions and the zeta-function of a singularity},
Comment. Math. Helv. \textbf{72} (1997), no.~2, 244--256.

\bibitem{LPS26}
M.~Larson, S.~Payne, and A.~Stapledon,
\emph{The local motivic monodromy conjecture for simplicial
nondegenerate singularities},
\href{https://arxiv.org/abs/2209.03553v3}{arXiv:2209.03553v3}
(2026).

\bibitem{LVP11}
A.~Lemahieu and L.~Van Proeyen,
\emph{Monodromy conjecture for non-degenerate surface singularities},
Trans. Amer. Math. Soc. \textbf{363} (2011), no.~9, 4801--4829.

\bibitem{Mil68}
J.~Milnor,
\emph{Singular Points of Complex Hypersurfaces},
Annals of Mathematics Studies, vol.~61,
Princeton University Press, Princeton, NJ, 1968.

\bibitem{Veys24}
W.~Veys,
\emph{Introduction to the monodromy conjecture},
\href{https://arxiv.org/abs/2403.03343v1}{arXiv:2403.03343v1}
(2024).

\end{thebibliography}
\end{document}
